\documentclass[11pt]{article}
\usepackage[margin=1in]{geometry}
\usepackage{amsmath,amssymb,amsthm,mathtools}
\usepackage{array,booktabs}
\usepackage{graphicx}
\usepackage{xcolor}
\usepackage[hidelinks]{hyperref}
\usepackage{microtype}

\definecolor{accent}{RGB}{37,86,123}
\newtheorem{theorem}{Theorem}
\newtheorem{corollary}[theorem]{Corollary}
\newtheorem{lemma}[theorem]{Lemma}
\theoremstyle{definition}
\newtheorem{definition}[theorem]{Definition}
\newtheorem{example}[theorem]{Example}

\newcommand{\Vol}{\operatorname{Vol}}
\newcommand{\R}{\mathbb{R}}

\title{\textcolor{accent}{Segment-sum concavity for projectively unimodular zonotopes}}
\author{Automated mathematics run \texttt{fa\_banach\_001}, lane 13}
\date{23 July 2026}

\begin{document}
\maketitle

\begin{center}
\fbox{\parbox{0.88\textwidth}{
\textbf{Status.} Substantial partial result; current verdict: likely valid.
The source conjecture remains open for general zonoids.
}}
\end{center}

\section{Source problem}

Fradelizi, Madiman, Meyer, and Zvavitch ask the following in Section~5 of
\cite{FMMZ}.  If \(Z\) is a zonoid in \(\R^n\),
\[
 L=\sum_{j=1}^{m}[0,u_j],\qquad m\leq n,
 \qquad
 P(t)=\Vol_n(Z+tL),
\]
is \(P(t)^{1/m}\) concave for \(t\geq0\)?  Their exact passage appears below.

\begin{center}
\includegraphics[width=0.96\textwidth]{figures/open_problem_crop.png}
\end{center}

The source proves nearby low-dimensional and two-segment results.  The theorem
below gives the conjectured conclusion in every dimension for a natural
regular-matroid class of finite zonotopes.

\section{Definitions}

\begin{definition}
A list of nonzero vectors \(w_1,\ldots,w_N\in\R^n\) is
\emph{projectively unimodular} if there are an invertible
\(T\in\operatorname{GL}_n(\R)\), positive scalars \(c_1,\ldots,c_N\), and the
columns \(a_1,\ldots,a_N\) of a full-row-rank totally unimodular
\(n\times N\) matrix \(A\) such that
\[
 w_e=c_eTa_e\qquad(1\leq e\leq N).
\]
\end{definition}

Replacing a generator by its negative only translates the corresponding
segment, so it leaves all relevant volumes unchanged; signs may therefore be
absorbed into the columns of \(A\).  Reduced oriented incidence matrices of
graphs are totally unimodular, hence graphical zonotopes supply a large
family of examples.

\section{Proof intuition}

For a general zonotope, its volume is a sum of absolute maximal determinants.
Absolute values normally prevent that sum from being a single determinant.
Projective unimodularity removes precisely this obstruction: every maximal
minor of \(A\) is \(0\) or \(\pm1\), so absolute minors equal squared minors.
Cauchy--Binet then packages the entire volume polynomial into a
positive-semidefinite determinant.

After the base generators receive weight \(1\) and the \(m\) distinguished
generators receive weight \(t\), the polynomial is \(\det(B+tC)\) with
\(B>0\) and \(C\geq0\).  Such a polynomial has only negative real roots.
The \(m\)-th root of a product of at most \(m\) positive affine factors is
concave by Cauchy--Schwarz.  This last observation also handles the case in
which the degree is smaller than \(m\).

\section{Main theorem}

\begin{theorem}\label{thm:main}
Let
\[
 Z=\sum_{i=1}^{r}[0,v_i]\subset\R^n
\]
be full-dimensional, let
\[
 L=\sum_{j=1}^{m}[0,u_j],\qquad 1\leq m\leq n,
\]
and suppose that the combined list
\((v_1,\ldots,v_r,u_1,\ldots,u_m)\) is projectively unimodular.
Then
\[
 P(t)=\Vol_n(Z+tL)
\]
has only negative real zeros.  Consequently \(P(t)^{1/m}\) is concave on
\([0,\infty)\).
\end{theorem}

\begin{proof}
Write \(N=r+m\), let \(w_1,\ldots,w_N\) denote the combined list, and choose
\(T,c_e,A=(a_1\ \cdots\ a_N)\) as in the definition.  The standard zonotope
volume formula gives, for nonnegative weights \(x_e\),
\begin{align}
\Vol_n\left(\sum_{e=1}^{N}x_e[0,w_e]\right)
 &=\sum_{\substack{S\subseteq[N]\\|S|=n}}
   |\det(w_e:e\in S)|\prod_{e\in S}x_e \notag\\
 &=|\det T|
   \sum_{\substack{S\subseteq[N]\\|S|=n}}
   |\det A_S|\prod_{e\in S}c_ex_e. \label{eq:zonotope}
\end{align}
Total unimodularity says \(\det A_S\in\{0,\pm1\}\), and hence
\(|\det A_S|=(\det A_S)^2\).  Cauchy--Binet turns \eqref{eq:zonotope} into
\begin{equation}\label{eq:determinant}
\Vol_n\left(\sum_{e=1}^{N}x_e[0,w_e]\right)
=|\det T|\det\!\left(
A\operatorname{diag}(c_ex_e)A^{\mathsf T}\right).
\end{equation}

Partition \(A=(A_Z\ A_L)\) and the \(c_e\)'s according to the generators of
\(Z\) and \(L\).  Setting the \(Z\)-weights to \(1\) and the \(L\)-weights to
\(t\) in \eqref{eq:determinant} yields
\[
P(t)=|\det T|\det(B+tC),
\]
where
\[
B=A_Z\operatorname{diag}(c_1,\ldots,c_r)A_Z^{\mathsf T},
\qquad
C=A_L\operatorname{diag}(c_{r+1},\ldots,c_{r+m})A_L^{\mathsf T}.
\]
Both matrices are positive semidefinite.  Since \(Z\) is full-dimensional,
the vectors \(v_1,\ldots,v_r\), and therefore the columns of \(A_Z\), span
\(\R^n\).  Thus \(B\) is positive definite.

If \(\operatorname{Im}t>0\) and \((B+tC)q=0\) for a complex vector \(q\), then
taking the imaginary part of \(q^*(B+tC)q=0\) gives
\[
(\operatorname{Im}t)\,q^*Cq=0.
\]
Hence \(Cq=0\), and then \(q^*Bq=0\), which forces \(q=0\).  Therefore
\(\det(B+tC)\neq0\) in the open upper half-plane.  Because \(P\) has real
coefficients, any nonreal zero would bring its conjugate into that half-plane.
Thus all zeros of \(P\) are real.  Its coefficients are nonnegative and
\(P(0)=\Vol_n(Z)>0\), so all its zeros are strictly negative.

Let \(d=\deg P\leq m\).  There are \(\rho_1,\ldots,\rho_d>0\) and \(K>0\)
such that
\[
P(t)=K\prod_{\ell=1}^{d}(t+\rho_\ell).
\]
For \(F=P^{1/m}\) and \(t\geq0\),
\[
\frac{F''(t)}{F(t)}
=-\frac1m\sum_{\ell=1}^{d}\frac1{(t+\rho_\ell)^2}
+ \frac1{m^2}\left(\sum_{\ell=1}^{d}\frac1{t+\rho_\ell}\right)^2.
\]
Cauchy--Schwarz gives
\[
\left(\sum_{\ell=1}^{d}\frac1{t+\rho_\ell}\right)^2
\leq d\sum_{\ell=1}^{d}\frac1{(t+\rho_\ell)^2}
\leq m\sum_{\ell=1}^{d}\frac1{(t+\rho_\ell)^2}.
\]
It follows that \(F''(t)\leq0\), as claimed.
\end{proof}

\begin{corollary}[Graphical zonotopes]\label{cor:graphical}
Let \(G\) be a connected graph on \(n+1\) vertices and represent each edge by
its column in a reduced oriented incidence matrix.  Partition the edges as
\(E=E_Z\sqcup E_L\), where the subgraph \((V,E_Z)\) is connected and
\(|E_L|=m\leq n\).  Positive rescalings of the edge vectors are allowed.
If \(Z\) is generated by \(E_Z\) and \(L\) by \(E_L\), then
\(\Vol_n(Z+tL)^{1/m}\) is concave for \(t\geq0\).
\end{corollary}

\begin{proof}
A reduced oriented incidence matrix is totally unimodular.  Connectivity of
\((V,E_Z)\) gives full rank to the \(Z\)-columns.  Apply
Theorem~\ref{thm:main}.
\end{proof}

\section{Exact four-dimensional example}

\begin{example}
Take the reduced incidence matrix of \(K_5\), so the ambient dimension is
\(4\).  Let the three distinguished edges be
\[
E_L=\{\{0,1\},\{0,2\},\{0,3\}\}
\]
and let the other seven edges generate \(Z\).  The latter subgraph is
connected, and exact Cauchy--Binet evaluation gives
\[
P(t)=4t^3+33t^2+72t+16=(t+4)^2(4t+1).
\]
This is a non-parallelotope example in the first dimension not covered by the
source paper's \(\R^3\) result.  Directly,
\[
3P(t)P''(t)-2P'(t)^2=-450(t+4)^2\leq0,
\]
which is the differential inequality for concavity of \(P^{1/3}\).
The exact verifier included with this packet checks all of these identities
over the integers.
\end{example}

\section{Verification and limitations}

\subsection*{Adversarial checks}

\begin{center}
\begin{tabular}{@{}>{\raggedright\arraybackslash}p{0.28\textwidth}
>{\raggedright\arraybackslash}p{0.63\textwidth}@{}}
\toprule
Potential failure & Check \\
\midrule
Missing determinant normalization &
For segments \([0,w_e]\), the zonotope formula is exactly the sum of absolute
maximal determinants; there is no factorial. \\
Absolute versus squared minors &
The equality uses total unimodularity explicitly.  It is false for a general
vector configuration. \\
Singular base pencil &
Full-dimensionality of \(Z\) is used to prove \(B>0\). \\
Degree smaller than \(m\) &
The final Cauchy--Schwarz bound uses \(d\leq m\), not \(d=m\). \\
Complex-kernel argument &
\(B,C\) are real symmetric positive semidefinite, so \(q^*Cq=0\) implies
\(Cq=0\), and \(q^*Bq=0\) implies \(q=0\). \\
\bottomrule
\end{tabular}
\end{center}

\subsection*{Scope}

The theorem does not settle the source conjecture for arbitrary zonoids, or
even arbitrary finite zonotopes.  Its hypothesis concerns the combined
directions of \(Z\) and \(L\).  A tempting attempt to divide a forced monomial
from the general degree-\(n\) Lorentzian volume polynomial fails because it
does not preserve the stronger degree-\(m\) Newton normalization.

\subsection*{Bounded novelty check}

On 23 July 2026, searches were run for the exact source conjecture and arXiv
id, and for combinations of \emph{zonoid}, \emph{totally unimodular},
\emph{projectively unimodular}, \emph{graphical zonotope},
\emph{volume polynomial}, \emph{real stable}, and \emph{concavity}.
The search recovered the source paper and general literature on Lorentzian
and volume polynomials, including \cite{BH}, but no statement of the precise
subcase above.  This is a bounded negative search, so novelty confidence is
moderate rather than definitive.

\subsection*{Human-review recommendation}

Send to a human reviewer.  The key steps to inspect are
\eqref{eq:determinant}, the use of full-dimensionality to obtain \(B>0\), and
the exact scope of the projectively-unimodular hypothesis.  If accepted, the
result is suitable as a short proposition or note; the general conjecture
should continue to be presented as open.

\begin{thebibliography}{9}

\bibitem{FMMZ}
M. Fradelizi, M. Madiman, M. Meyer, and A. Zvavitch,
\emph{On the volume of the Minkowski sum of zonoids},
J. Funct. Anal. \textbf{286} (2024), 110247;
\href{https://arxiv.org/abs/2206.02123}{arXiv:2206.02123}.

\bibitem{BH}
P. Br\"and\'en and J. Huh,
\emph{Lorentzian polynomials},
Ann. of Math. \textbf{192} (2020), 821--891;
\href{https://arxiv.org/abs/1902.03719}{arXiv:1902.03719}.

\end{thebibliography}

\end{document}
